\chapter[1931 --- Warburg]{1931: Otto Heinrich Warburg}
\label{ch:1931_ottoheinrichwarburg}
\index{Warburg, Otto Heinrich|textbf}
\index{Atmungsferment|textbf}
\index{Cellular respiration|textbf}
\index{Cytochrome oxidase|textbf}
\index{Warburg effect|textbf}
\index{Aerobic glycolysis|textbf}
\index{Warburg manometer|textbf}
\index{FDG-PET imaging|textbf}

\section*{Main Contribution}

\textit{Discovered the nature and mode of action of the iron-containing respiratory enzyme (Atmungsferment) that catalyzes cellular biological oxidation, and introduced quantitative manometry to the study of tissue metabolism and tumor bioenergetics.}\cite{nobel-1931,lecture-1931,warburg-cancer-review}

\section{Official Citation}

for his discovery of the nature and mode of action of the respiratory enzyme

\section{Background and Historical Context}

At the start of the twentieth century, the physical-chemical nature of biological combustion remained one of biology's most contentious mysteries. While Antoine Lavoisier had demonstrated in the late eighteenth century that respiration is a slow combustion process consuming oxygen and generating carbon dioxide and heat, how this occurred within the watery, neutral-pH environment of living cells was unexplained. 

Two opposing theories divided biochemists. The English school, led by Heinrich Wieland, argued for ``dehydrogenation'': organic substrates were supposedly activated by dehydrases that stripped hydrogen atoms from fuel molecules, transferring them spontaneously to molecular oxygen. In contrast, the German school argued for ``oxygen activation.'' Otto Heinrich Warburg, trained in chemistry under Emil Fischer in Berlin and in medicine in Heidelberg, insisted that living cells contain specific catalytic enzymes that directly activate molecular oxygen.

Warburg, working at the Kaiser Wilhelm Institute for Biology in Berlin-Dahlem, realized that settling this dispute required new quantitative instrumentation to measure minute rates of oxygen consumption in living tissue slices under physiological conditions.

\section{The Discovery of the Respiratory Enzyme (Atmungsferment)}

\subsection{The Warburg Manometer and Tissue Slice Technique}

Warburg's first great methodological triumph was the invention of the \emph{Warburg manometer} and the development of the thin-tissue-slice method. Before Warburg, tissue homogenates rapidly lost their respiration. Warburg showed that by cutting razor-thin tissue sections (less than 0.3 millimeters thick) and suspending them in buffered saline solutions in a temperature-controlled water bath, oxygen could diffuse freely into all cells without encountering hypoxia in the center. 

By measuring the microscopic displacement of manometric fluid in a calibrated U-tube as oxygen was absorbed from the gas space above the tissue, Warburg could quantify cellular respiration ($\text{Q}_{\text{O}_2}$, microliters of oxygen consumed per milligram dry weight per hour) with unprecedented accuracy.

\subsection{Probing the Respiratory Enzyme with Inhibitors and Light}

Warburg observed that trace amounts of cyanide ($\text{CN}^-$) or carbon monoxide ($\text{CO}$) completely and reversibly paralyzed cellular respiration in living yeast, sea urchin eggs, and mammalian tissues. Because cyanide and carbon monoxide were known to bind transition metals—particularly iron in hemoproteins—Warburg deduced that the principal respiratory catalyst in cells must be an iron-containing porphyrin complex, which he termed the \emph{Atmungsferment} (respiratory enzyme).

To prove that the enzyme was a hemochromogen without being able to isolate or crystallize the membrane-bound protein, Warburg devised a brilliant photochemical experiment. He exploited the known fact that carbon monoxide-iron complexes are photolabile: illuminating the system with bright light dissociates the $\text{Fe-CO}$ bond, restoring the enzyme's catalytic activity. 

Warburg placed carbon monoxide-poisoned, non-respiring yeast cells in the dark and irradiated them with monochromatic light across different wavelengths of the spectrum, measuring the precise degree to which cellular respiration was restored at each wavelength. Because the photochemical efficiency of dissociation is proportional to the light absorbed by the $\text{Fe-CO}$ complex, plotting respiration restoration against wavelength produced the exact \emph{photochemical absorption spectrum} of the respiratory enzyme. 

The resulting action spectrum revealed distinct absorption bands in the ultraviolet (Soret band around 430 nm) and visible regions (around 590 nm), precisely matching the absorption spectrum of an iron-porphyrin hemochrome similar to hemoglobin and myoglobin. This definitive demonstration proved that cellular respiration is mediated by an iron-containing porphyrin enzyme, now known as \textbf{cytochrome c oxidase} (Complex IV of the mitochondrial electron transport chain). For this discovery, Warburg was awarded the 1931 Nobel Prize in Physiology or Medicine.

\section{Tumor Metabolism and the Warburg Effect}

In parallel with his studies on normal respiration, Warburg applied his manometric methods to malignant tumor tissues. In the 1920s, he made an observation that would influence cancer biology for a century: whereas normal differentiated adult tissues use glycolysis to produce pyruvate only in the absence of oxygen (fermentation), switching to mitochondrial oxidative phosphorylation in the presence of oxygen (the Pasteur effect), malignant tumor cells behave fundamentally differently.

Cancer cells consume enormous quantities of glucose—often ten to fifty times more than their non-malignant tissue counterparts—and convert the vast majority into lactic acid, \emph{even in the presence of abundant oxygen}. This phenomenon is universally known as **aerobic glycolysis** or the **Warburg effect**.

Warburg later hypothesized that the irreversible injury of cellular respiration, forcing cells to revert to primitive anaerobic fermentation, was the universal primary cause of cancer. While Warburg's causal hypothesis (that respiration impairment initiates cancer) was subsequently disproved—modern oncology established that somatic mutations in oncogenes and tumor suppressors drive malignant transformation, and many tumors retain functional, actively respiring mitochondria—his physiological observation remains one of the hallmark metabolic features of cancer.

\section{Modern Molecular Mechanisms of the Warburg Effect}

Contemporary biochemistry has revealed why cancer cells adopt the Warburg phenotype:
\begin{itemize}
    \item \textbf{Oncogenic Reprogramming}: Mutations in oncogenes such as *MYC*, *KRAS*, and the *PI3K/Akt/mTOR* pathway directly upregulate glucose transporters (specifically GLUT1 and GLUT3) and rate-limiting glycolytic enzymes such as hexokinase 2 (HK2) and pyruvate kinase M2 (PKM2).
    \item \textbf{Hypoxia-Inducible Factor 1 (HIF-1)}: HIF-1$\alpha$ activation (whether through tumor hypoxia or pseudohypoxia caused by *VHL* mutations or succinate/fumarate accumulation) drives expression of lactate dehydrogenase A (LDHA) and pyruvate dehydrogenase kinase 1 (PDK1), which phosphorylates and inactivates pyruvate dehydrogenase, shunting pyruvate away from mitochondria and into lactate.
    \item \textbf{Biosynthetic Precursor Generation}: Rapidly dividing cancer cells require more than just ATP; they require massive quantities of building blocks—ribose-5-phosphate for nucleotides, glycerol and fatty acids for lipid membranes, and non-essential amino acids for protein synthesis. Diverting glycolytic intermediates into the pentose phosphate pathway and one-carbon metabolism supports rapid macromolecular biosynthesis far more efficiently than burning every glucose molecule completely to $\text{CO}_2$ in mitochondria.
\end{itemize}

\section{Impact on Modern Medicine and Clinical Oncology}

Warburg's discoveries exert an immense influence across contemporary clinical practice:
\begin{itemize}
    \item \textbf{Positron Emission Tomography (FDG-PET)}: The single most impactful clinical translation of the Warburg effect is 2-deoxy-2-[$^{18}\text{F}$]fluoro-D-glucose (FDG) PET imaging. Tumor cells greedily import the radiolabeled glucose analog via overexpressed GLUT transporters, where it is phosphorylated by hexokinase into FDG-6-phosphate and trapped intracellularly. FDG-PET is a cornerstone of clinical oncology worldwide for staging tumors, assessing therapeutic responses, and detecting cancer recurrence across lymphomas, lung cancer, colorectal cancer, and melanoma.
    \item \textbf{Cancer Metabolism Therapeutics}: Small-molecule inhibitors targeting cancer metabolic vulnerabilities—including lactate dehydrogenase (LDHA), monocarboxylate transporters (MCT1/4 inhibitors to prevent lactate extrusion and acidify tumors), glutaminase (GLS) inhibitors, and isocitrate dehydrogenase 1/2 (IDH1/IDH2) inhibitors (such as ivosidenib and enasidenib for acute myeloid leukemia and gliomas)—are direct descendants of Warburg's metabolic paradigm.
    \item \textbf{Mitochondrial Medicine}: Understanding cytochrome c oxidase, iron-sulfur clusters, and electron transport kinetics underpins the diagnosis and molecular management of inherited mitochondrial respiratory chain disorders (such as Leigh syndrome and MELAS).
\end{itemize}

\section{Legacy}

Otto Heinrich Warburg stands as one of the titans of twentieth-century biochemistry. Over a career spanning six decades, he trained or mentored a constellation of future Nobel laureates, including Hans Krebs (Krebs cycle), Fritz Lipmann (coenzyme A), and George Wald (biochemistry of vision). His unwavering insistence on rigorous, direct quantitative physical measurements permanently transformed physiology into a molecular science.

\section{Bibliography}

\begin{thebibliography}{9}
\bibitem{nobel-1931}
Nobel Prize Outreach, ``The Nobel Prize in Physiology or Medicine 1931,''\emph{NobelPrize.org}.\\
\url{https://www.nobelprize.org/prizes/medicine/1931/summary/}

\bibitem{lecture-1931}
Otto Heinrich Warburg, ``The Oxygen-Transferring Ferment of Respiration,'' Nobel Lecture, December 10, 1931.\\
\url{https://www.nobelprize.org/prizes/medicine/1931/warburg/lecture/}

\bibitem{warburg-classic-1926}
O. Warburg, \emph{Über den Stoffwechsel der Tumoren} (Berlin: Springer, 1926); English translation by F. Dickens, \emph{The Metabolism of Tumours} (London: Constable \& Co., 1930).

\bibitem{warburg-science-1956}
O. Warburg, ``On the Origin of Cancer Cells,''\emph{Science} 123 (1956): 309--314.\\
\url{https://doi.org/10.1126/science.123.3191.309}

\bibitem{warburg-cancer-review}
M. G. Vander Heiden, L. C. Cantley, and C. B. Thompson, ``Understanding the Warburg Effect: The Metabolic Requirements of Cell Proliferation,''\emph{Science} 324 (2009): 1029--1033.\\
\url{https://doi.org/10.1126/science.1160809}

\bibitem{hay-warburg-revisited}
N. Hay, ``Reprogramming glucose metabolism in cancer: can it be exploited for cancer therapy?,''\emph{Nature Reviews Cancer} 16 (2016): 635--649.\\
\url{https://doi.org/10.1038/nrc.2016.77}
\end{thebibliography}
